% test-012.tex -- comandos \spcoord, \spPoint y \spvec
%
% Compilar:  lualatex-dev test-012.tex
% Validar:   show-pdf-tags --xml test-012.pdf | rnv-wrapp
%            verapdf --flavour ua2 --format text test-012.pdf
%
% Cada caso es un parrafo numerado, para poder referirse a el al escuchar
% el PDF. Los comentarios "doc:" indican lo que documenta el manual; en el
% resto, la lectura se revisa escuchando. Todos trabajan en modo
% matematico. Reglas de la familia geo: un punto es una unica letra
% mayuscula con prima o con subindice entre llaves (A_{1}, nunca A_1); un
% vector con nombre es una letra minuscula. Las coordenadas (x, y) usan
% coma, o punto y coma si llevan coma decimal. Las entradas invalidas (que
% detienen la compilacion) no van aqui, entre ellas A_1, un punto en
% minuscula, \spcoord con uno o con tres valores y \spvec{ab}.
\DocumentMetadata
  {
    lang=es-CL,
    pdfversion=2.0,
    pdfstandard=ua-2,
    tagging=on,
    tagging-setup={math/setup={mathml-SE}},
  }
\documentclass{article}
\usepackage[spanish]{babel}
\usepackage{unicode-math}
\usepackage{spintent}
\usepackage{hyperref}
\hypersetup
  {
    pdfdisplaydoctitle = true,
    pdftitle           = {test-012: spcoord, spPoint y spvec},
  }
\pagestyle{empty}
\begin{document}

\section*{A. Comando spcoord}

% doc: «menos 3 coma 4»
1. Enteros con signo: $\spcoord{-3, 4}$.

% doc: «x subíndice 1 coma y subíndice 1»
2. Con subíndices: $\spcoord{x_{1},y_{1}}$.

% doc: «3 coma 4»
3. Enteros positivos: $\spcoord{3, 4}$.

% doc: «0 coma 0»
4. El origen: $\spcoord{0, 0}$.

% doc: «menos 3 coma 4»
5. Sin espacio tras la coma: $\spcoord{-3,4}$.

% doc: «3 punto 5 coma 4 punto 2»
6. Decimales con punto: $\spcoord{3.5, 4.2}$.

% doc: «3 coma 5 punto y coma 4 coma 2»
7. Decimales con coma y punto y coma: $\spcoord{3,5; 4,2}$.

% doc: «menos 1 coma 5 punto y coma 2»
8. Con signo y punto y coma: $\spcoord{-1,5; 2}$.

% doc: «x coma y»
9. Dos letras: $\spcoord{x, y}$.

% doc: «a coma b»
10. Otras letras: $\spcoord{a, b}$.

% doc: «a más 1 coma b menos 1»
11. Con expresiones: $\spcoord{a+1, b-1}$.

% doc: «1 barra oblicua 2 coma 3 barra oblicua 4»
12. Con cocientes en línea: $\spcoord{1/2, 3/4}$.

% doc: «1 mitad coma 3»
13. Con una fracción: $\spcoord{\frac{1}{2}, 3}$.

% doc: «el raíz cuadrada de 2 coma 1»
14. Con una raíz: $\spcoord{\sqrt{2}, 1}$.

% doc: «pi coma 2»
15. Con pi: $\spcoord{\pi, 2}$.

% doc: «2 pi coma 0»
16. Con un múltiplo de pi: $\spcoord{2\pi, 0}$.

% doc: «menos 3 coma 4»
17. Con silent false, el valor por defecto: $\spcoord[silent=false]{-3, 4}$.

% doc: «»
18. Con silent true: $\spcoord[silent=true]{-3, 4}$.

% doc: «»
19. Con silent true y letras: $\spcoord[silent=true]{x_{1},y_{1}}$.

\section*{B. Comando spPoint: el punto}

% doc: «punto A»
20. Una letra: $\spPoint{A}$.

% doc: «punto A sub uno»
21. Con subíndice: $\spPoint{A_{1}}$.

% doc: «punto A prima»
22. Con una prima: $\spPoint{A'}$.

% doc: «punto A doble prima»
23. Con dos primas: $\spPoint{A''}$.

% doc: «punto A sub 12»
24. Con subíndice de dos cifras: $\spPoint{A_{12}}$.

% doc: «punto A sub n»
25. Con subíndice literal: $\spPoint{A_{n}}$.

% doc: «punto B»
26. Otra letra: $\spPoint{B}$.

% doc: «punto Z»
27. La última letra: $\spPoint{Z}$.

% doc: «punto A sub uno punto B»
28. Dos puntos seguidos, el primero con subíndice: $\spPoint{A_{1}}\spPoint{B}$.

% doc: «punto A coma punto B»
29. Dos puntos seguidos con coma: $\spPoint{A}, \spPoint{B}$.

\section*{C. Comando spPoint: con coordenadas}

% doc: «punto A menos 3 coma 4»
30. Con coordenadas: $\spPoint{A}(-3,4)$.

% doc: «punto A sub uno menos 3 coma 4»
31. Con subíndice y coordenadas: $\spPoint{A_{1}}(-3,4)$.

% doc: «punto A prima menos 3 coma 4»
32. Con prima y coordenadas: $\spPoint{A'}(-3,4)$.

% doc: «punto A 3 punto 5 coma 4 punto 2»
33. Con decimales con punto: $\spPoint{A}(3.5, 4.2)$.

% doc: «punto A 3 coma 5 punto y coma 4 coma 2»
34. Con decimales con coma y punto y coma: $\spPoint{A}(3,5; 4,2)$.

% doc: «punto P x coma y»
35. Con coordenadas literales: $\spPoint{P}(x, y)$.

% doc: «punto O 0 coma 0»
36. El origen: $\spPoint{O}(0, 0)$.

\section*{D. Comando spPoint: llaves}

% doc: «punto A»
37. Con read-cmd school, el valor por defecto: $\spPoint[read-cmd=school]{A}$.

% doc: «punto a»
38. Con read-cmd mathcat: $\spPoint[read-cmd=mathcat]{A}$.

% doc: «punto a subíndice 1»
39. Con read-cmd mathcat y subíndice: $\spPoint[read-cmd=mathcat]{A_{1}}$.

% doc: «punto A»
40. Con concept true, el valor por defecto: $\spPoint[concept=true]{A}$.

% doc: «A»
41. Con concept false: $\spPoint[concept=false]{A}$.

% doc: «A sub uno»
42. Con concept false y subíndice: $\spPoint[concept=false]{A_{1}}$.

% doc: «A menos 3 coma 4»
43. Con concept false y coordenadas: $\spPoint[concept=false]{A}(-3,4)$.

% doc: «el punto A»
44. Con prefix: $\spPoint[prefix={el}]{A}$.

% doc: «el punto A menos 3 coma 4»
45. Con prefix y coordenadas: $\spPoint[prefix={el}]{A}(-3,4)$.

% doc: «punto A»
46. Con silent false, el valor por defecto: $\spPoint[silent=false]{A}$.

% doc: «»
47. Con silent true: $\spPoint[silent=true]{A}$.

% doc: «»
48. Con silent true y coordenadas: $\spPoint[silent=true]{A_{1}}(-3,4)$.

% doc: «a»
49. Con read-cmd mathcat y concept false: $\spPoint[read-cmd=mathcat,concept=false]{A}$.

% doc: «el punto a subíndice 1 3 coma 5 punto y coma 4 coma 2»
50. Con todas las llaves: $\spPoint[read-cmd=mathcat,prefix={el},concept=true,silent=false]{A_{1}}(3,5; 4,2)$.

\section*{E. Comando spvec: vector con nombre}

% doc: «vector a»
51. Una letra: $\spvec{a}$.

% doc: «vector a sub uno»
52. Con subíndice: $\spvec{a_{1}}$.

% doc: «vector u»
53. Otra letra: $\spvec{u}$.

% doc: «vector a es igual a menos 3 coma 4»
54. Con coordenadas: $\spvec{a}(-3,4)$.

% doc: «vector a sub uno es igual a 2 coma 5»
55. Con subíndice y coordenadas: $\spvec{a_{1}}(2,5)$.

% doc: «vector v es igual a 3 coma 5 punto y coma 4 coma 2»
56. Con decimales y coordenadas: $\spvec{v}(3,5; 4,2)$.

\section*{F. Comando spvec: vector definido por dos puntos}

% doc: «vector A B»
57. Dos puntos seguidos: $\spvec{AB}$.

% doc: «vector A B»
58. Dos puntos con coma: $\spvec{A,B}$.

% doc: «vector A B»
59. Dos puntos con coma y espacio: $\spvec{A, B}$.

% doc: «vector A prima B prima»
60. Con primas: $\spvec{A'B'}$.

% doc: «vector A sub uno B sub dos»
61. Con subíndices: $\spvec{A_{1}B_{2}}$.

% doc: «vector A sub uno B sub dos»
62. Con subíndices y coma: $\spvec{A_{1},B_{2}}$.

% doc: «vector A B es igual a menos 3 coma 4»
63. Con coordenadas: $\spvec{AB}(-3,4)$.

% doc: «vector M N es igual a x coma y»
64. Con coordenadas literales: $\spvec{MN}(x, y)$.

\section*{G. Comando spvec: llaves}

% doc: «vector A B»
65. Con read-cmd school, el valor por defecto: $\spvec[read-cmd=school]{AB}$.

% doc: «vector a b»
66. Con read-cmd mathcat: $\spvec[read-cmd=mathcat]{AB}$.

% doc: «el vector a»
67. Con prefix y un nombre: $\spvec[prefix={el}]{a}$.

% doc: «el vector A B»
68. Con prefix y dos puntos: $\spvec[prefix={el}]{AB}$.

% doc: «el vector a es igual a menos 3 coma 4»
69. Con prefix y coordenadas: $\spvec[prefix={el}]{a}(-3,4)$.

% doc: «vector a»
70. Con autofit false, el valor por defecto: $\spvec[autofit=false]{a}$.

% doc: «vector a»
71. Con autofit true: $\spvec[autofit=true]{a}$.

% doc: «vector a»
72. Con concept true, el valor por defecto: $\spvec[concept=true]{a}$.

% doc: «a»
73. Con concept false y un nombre: $\spvec[concept=false]{a}$.

% doc: «A B»
74. Con concept false y dos puntos: $\spvec[concept=false]{AB}$.

% doc: «A B es igual a menos 3 coma 4»
75. Con concept false y coordenadas: $\spvec[concept=false]{AB}(-3,4)$.

% doc: «vector a»
76. Con silent false, el valor por defecto: $\spvec[silent=false]{a}$.

% doc: «»
77. Con silent true y un nombre: $\spvec[silent=true]{a}$.

% doc: «»
78. Con silent true y dos puntos: $\spvec[silent=true]{AB}$.

% doc: «»
79. Con silent true y coordenadas: $\spvec[silent=true]{a}(-3,4)$.

% doc: «el vector a subíndice 1 es igual a 3 coma 5 punto y coma 4 coma 2»
80. Con todas las llaves: $\spvec[read-cmd=mathcat,prefix={el},autofit=true,concept=true,silent=false]{a_{1}}(3,5; 4,2)$.

\section*{H. En fórmulas}

% doc: «vector a más vector b es igual a vector c»
81. Suma de vectores: $\spvec{a} + \spvec{b} = \spvec{c}$.

% doc: «vector A B es igual a vector C D»
82. Igualdad de vectores: $\spvec{AB} = \spvec{CD}$.

% doc: «vector a es igual a 1 coma 2, vector b es igual a 3 coma 4 y vector c es igual a 4 coma 6»
83. Tres vectores con coordenadas, cada uno en su fórmula: $\spvec{a}(1,2)$, $\spvec{b}(3,4)$ y $\spvec{c}(4,6)$.

% doc: «punto A 1 coma 2 vector A B es igual a 3 coma 4»
84. Un punto y un vector: $\spPoint{A}(1,2) \quad \spvec{AB}(3,4)$.

% doc: «punto A es igual a 1 coma 2»
85. Punto y coordenadas por separado: $\spPoint{A} = \spcoord{1, 2}$.

% doc: «vector A B es igual a punto B 1 coma 2 menos punto A 4 coma menos 2»
86. En fórmula destacada:
\[ \spvec{AB} = \spPoint{B}(1,2) - \spPoint{A}(4,-2) \].

\end{document}
